================================================================
  FINAL DISSERTATION REPORT — SAHIL SAXENA (11409565)
  A Human-in-the-Loop Framework for Alzheimer's Disease
  Stage Classification Using Multimodal Biomarkers
  University of Manchester, 2025–2026

  This file contains all written sections produced/revised
  in collaboration with Claude. Append new sections below.
================================================================


================================================================
  CHAPTER 3: TECHNICAL BACKGROUND AND LITERATURE REVIEW
  Sections 3.1 -- 3.4
================================================================

% -----------------------------------------------------------
\section{Machine Learning for Alzheimer's Disease Classification}
\label{sec:ml-ad}
% -----------------------------------------------------------

The application of supervised machine learning to Alzheimer's disease
classification has a history spanning two decades, evolving in tandem with both
algorithmic advances and the growing availability of multimodal biomarker
datasets. \citet{kloppel2008automatic} established an early proof of concept by
demonstrating that a linear support vector machine (SVM) operating on
voxel-based morphometry features derived from structural MRI could discriminate
AD from cognitively normal individuals with accuracy comparable to expert
radiologists. That result set the benchmark for automated biomarker-based
classification, but also exposed its primary limitation: the reliance on a
single imaging modality and a binary label, neither of which reflects the
clinical reality of three-class staging from heterogeneous tabular data.

The subsequent diversification of methods is systematically reviewed by
\citet{rathore2017review}, who document the consistent advantage of multi-modal
feature sets over single-modality inputs across more than 300 classification
studies. Critically for the present work, their review identifies three
structural challenges that apply directly to ADNI-derived cohorts: class
imbalance between diagnostic groups, small effective sample sizes relative to
feature dimensionality, and variable missingness arising from differential
availability of CSF, plasma, and imaging biomarkers across participants. Each of
these challenges shaped design decisions in this project: the near-balanced
576-patient cohort mitigates the first, repeated stratified cross-validation
addresses the second, and fold-aware $k$-NN imputation handles the third.

Deep learning approaches have produced impressive results where sufficient
imaging data are available. \citet{ding2019deep} trained a convolutional network
on $^{18}$F-FDG PET scans and demonstrated detection of AD pathology up to six
years before clinical diagnosis. However, the gains offered by deep learning are
contingent on large labelled datasets; for tabular cohorts of the scale used
here (576 patients, 12 selected features), well-regularised ensemble methods
are empirically competitive, and their interpretability properties are
substantially more tractable for clinical deployment \citep{rathore2017review}.

Multi-class staging (SCD vs.\ MCI vs.\ AD) presents a qualitatively harder
problem than binary detection. MCI is the most diagnostically heterogeneous
category---spanning amnestic and non-amnestic subtypes with differing
progression trajectories---and the SCD--MCI boundary is defined by a continuous
gradient rather than a discrete threshold. This structural ambiguity manifests
as suppressed per-class accuracy at the MCI boundary in classifiers that lack an
uncertainty-aware output mechanism, motivating the uncertainty quantification
and escalation framework described in Sections~\ref{sec:uq} and~\ref{sec:hitl}.

% -----------------------------------------------------------
\section{Ensemble Methods and Model Stacking}
\label{sec:ensemble}
% -----------------------------------------------------------

Ensemble learning reduces generalisation error by combining the outputs of
multiple base learners whose errors are partially independent.
\citet{breiman2001random} showed that the Random Forest---combining bagging with
random feature subsampling at each split---achieves robust generalisation on
tabular data by decorrelating tree predictions, making it naturally resistant to
overfitting in small-to-medium cohorts. The gradient boosting paradigm,
formalised by \citet{chen2016xgboost} in XGBoost through second-order gradient
approximations and L1/L2 regularisation, yields state-of-the-art performance on
structured prediction benchmarks. \citet{prokhorenkova2018catboost} extended
this with \emph{ordered boosting}, which estimates gradients on subsets that
exclude the current sample, mitigating the target leakage that standard boosting
can introduce when gradient statistics are computed over the full training fold.
In this project, CatBoost, Random Forest, XGBoost, a neural network, and an SVM
are included as base learners specifically because they represent distinct
inductive biases---variance reduction, sequential correction, ordered
statistics, non-linear activation, and margin maximisation---whose errors are
unlikely to co-occur on the same cases.

\emph{Stacked generalisation} \citep{wolpert1992stacked} trains a meta-learner
on the out-of-fold (OOF) probability outputs of the base models, learning
optimal combinations of base model estimates rather than raw features. Training
on OOF predictions is essential: a base model that memorises training data will
produce near-perfect in-sample outputs, and a meta-learner trained on those same
predictions would upweight the overfitting model, producing misleadingly
optimistic performance. Here, a logistic regression meta-learner is trained on a
15-dimensional vector (five models $\times$ three class probabilities) and
compared against simple probability averaging; the better-performing combination
is selected for downstream use.

Probability calibration is a prerequisite for a reliable escalation threshold.
A model that reports 0.90 confidence but is correct only 75\% of the time
renders any threshold $\tau$ meaningless. \citet{platt1999probabilistic}
proposed sigmoid calibration to map raw scores to calibrated probabilities
post-hoc; \citet{niculescu2005predicting} extended this comparison to isotonic
regression and showed the relative advantage of each method depends on
calibration set size, with isotonic regression requiring more data to avoid
overfitting the calibration mapping. Both methods are evaluated here using
Expected Calibration Error (ECE) as the selection criterion; the raw ensemble
probabilities achieve ECE~$= 0.042$ (Section~\ref{sec:baseline}), outperforming
both post-hoc calibrators on this cohort, and are retained for all downstream
use.

% -----------------------------------------------------------
\section{Human-in-the-Loop Machine Learning}
\label{sec:hitl}
% -----------------------------------------------------------

A Human-in-the-Loop (HITL) system integrates human judgement into the machine
learning pipeline in a principled, anticipatory way: the system identifies where
human input adds the most value, presents the relevant information efficiently,
and closes the loop by using feedback to update the model
\citep{monarch2021hitl}. The defining distinction from conventional decision
support is selectivity: rather than routing every AI output to a clinician for
review, the system escalates only a calibrated subset of cases where the model's
uncertainty is elevated or data quality is insufficient.
\citet{amershi2014power} characterise this as the core mechanism by which
interactive ML preserves automation throughput while recovering the accuracy
benefits of human oversight on the cases that most require it.

The case for HITL deployment in medical AI is well-supported.
\citet{topol2019highperformance} argues that AI and clinical expertise are
complementary rather than competitive: AI performs best on high-volume,
high-confidence cases while clinicians add the greatest marginal value on
complex, uncertain ones. \citet{rajpurkar2022ai} survey the clinical AI
landscape and conclude that most high-performing medical AI systems still
require structured human oversight to be safe and deployable---a conclusion that
applies with particular force to AD staging, where misclassification determines
DMT eligibility and false positives expose patients to immunotherapy side
effects, including amyloid-related imaging abnormalities.

In this project, the escalation policy takes a three-tier form.
\textbf{AI-Autonomous} cases are those where all uncertainty signals fall below
calibrated thresholds; the model prediction is accepted without review.
\textbf{AI-Assisted} cases trigger a clinician review informed by the AI
recommendation and SHAP feature attributions. \textbf{Clinician-Mandatory} cases
involve critical data absence or cross-modality biomarker inconsistencies that
place the case outside the model's reliable operating region regardless of
predicted confidence. This graduated structure concentrates the review budget on
cases where a change in outcome is most likely, the quantitative benefit of
which is demonstrated in Section~\ref{sec:hitl-performance}.

% -----------------------------------------------------------
\section{Uncertainty Quantification and Conformal Prediction}
\label{sec:uq}
% -----------------------------------------------------------

Meaningful uncertainty quantification is the foundation on which the escalation
engine rests. The standard scalar measures for a probabilistic classifier are
\emph{maximum predicted probability} (confidence), \emph{decision margin} (the
gap between top-1 and top-2 class probabilities), and \emph{predictive entropy}:
\begin{equation}
  H(\mathbf{p}) = -\sum_{k=1}^{K} p_k \log p_k,
  \label{eq:entropy}
\end{equation}
where $p_k$ is the predicted probability for class $k$ and $K = 3$ in this
project. Entropy reaches its maximum of $\log K$ when probability is
distributed uniformly across classes, and zero for a perfectly confident
prediction. These scalar measures are complementary---a low-confidence
prediction may still carry a large margin if probability mass concentrates on
two classes, while high entropy captures three-way ambiguity that neither
confidence nor margin alone detects---but none carries a formal statistical
guarantee about coverage of the true label.

\emph{Conformal prediction} \citep{shafer2008tutorial,
angelopoulos2023conformal} addresses this gap by constructing
\emph{prediction sets}: sets of candidate labels that are guaranteed to contain
the true label with pre-specified probability $1 - \alpha$, without assumptions
on the data distribution or model calibration quality. A conformal predictor
computes a nonconformity score for each calibration sample (here, the complement
of the predicted probability for the true class), estimates the
$(1 - \alpha)$-quantile $\hat{q}$ of the empirical score distribution with a
finite-sample correction, and includes class $k$ in the prediction set if its
predicted probability exceeds $1 - \hat{q}$.

The \emph{Adaptive Prediction Sets} (APS) variant \citep{angelopoulos2023conformal}
improves on standard conformal prediction for multi-class problems by forming
nonconformity scores from the sorted cumulative probability distribution,
producing sets that shrink on easy cases and expand on genuinely ambiguous ones.
This adaptive property has direct clinical utility: a set of $\{\text{SCD}\}$
at $\alpha = 0.05$ is a strong, actionable recommendation, while a set of
$\{\text{SCD}, \text{MCI}, \text{AD}\}$ is an unambiguous signal for mandatory
escalation. An $\alpha$-sweep over $\{0.05, 0.10, 0.15, 0.20\}$ yields
empirical coverages of 1.000 across all settings (Section~\ref{sec:conformal}),
with average prediction set sizes ranging from 2.34 to 2.61---a result that
reflects the high inherent uncertainty in this three-class cohort and the
conservative nature of the finite-sample CP guarantee.


================================================================
  CHAPTER 2: CLINICAL BACKGROUND
================================================================

\chapter{Clinical Background}
\label{chap:clinical-background}

% -----------------------------------------------------------
\section{Alzheimer's Disease Pathology}
\label{sec:ad-pathology}
% -----------------------------------------------------------

Alzheimer's disease (AD) is the most prevalent neurodegenerative disorder
worldwide, accounting for an estimated 60--70\% of the approximately 57 million
people living with dementia globally \citep{who2023dementia}. Its molecular
pathology is defined by two hallmark lesions: extracellular neuritic plaques
composed of misfolded amyloid-$\beta$ (A$\beta$) peptide, and intraneuronal
neurofibrillary tangles (NFTs) formed by hyperphosphorylated tau protein. The
\emph{amyloid cascade hypothesis} posits that aberrant cleavage of amyloid
precursor protein by $\beta$- and $\gamma$-secretases generates
A$\beta_{42}$ fragments whose aggregation triggers a downstream cascade of
neuroinflammation, synaptic dysfunction, tau hyperphosphorylation, and
progressive neuronal loss \citep{mckhann2011diagnosis}.

The temporal ordering of these pathological events is now well characterised.
\citet{jack2018nia} formalised this understanding in the NIA-AA
\emph{A/T/(N) biomarker framework}, which groups biomarkers into three
categories: \textbf{A} (amyloid pathology), measured by CSF
A$\beta_{42}$, the A$\beta_{42}$/A$\beta_{40}$ ratio, or amyloid PET;
\textbf{T} (tau pathology), measured by CSF phosphorylated tau (p-tau) or tau
PET; and \textbf{(N)} (neurodegeneration), measured by CSF total tau (t-tau),
FDG-PET hypometabolism, or structural MRI atrophy. Crucially, amyloid
deposition is detectable approximately 15--20 years before the onset of
clinical symptoms, establishing a substantial preclinical window for
biomarker-guided intervention and---as this project demonstrates---for
machine-learning systems that can flag at-risk individuals before irreversible
cognitive decline.

More recently, plasma-based biomarkers have made non-invasive screening
increasingly feasible \citep{zetterberg2021plasma}. Plasma neurofilament light
chain (NfL) reflects axonal injury across neurodegenerative conditions, while
plasma p-tau reliably distinguishes AD from non-AD dementias. The availability
of both CSF and plasma markers within the same cohort---as in the ADNI dataset
used here---enables the multi-modal feature space that underpins this project's
ensemble classifier and, as the SHAP analysis in
Section~\ref{sec:explainability} reveals, exposes which biomarker modality
contributes most to each diagnostic boundary.

% -----------------------------------------------------------
\section{The AD Continuum: SCD, MCI, and Dementia}
\label{sec:ad-continuum}
% -----------------------------------------------------------

Contemporary research frames Alzheimer's disease not as a discrete clinical
event but as a decades-long biological continuum that unfolds in clinically
distinguishable stages. The first recognisable stage, \emph{Subjective
Cognitive Decline} (SCD), is defined by self-reported perception of worsening
memory or cognitive function in the absence of detectable objective impairment
on standardised neuropsychological tests \citep{jessen2014scd}. Although SCD
individuals perform within normal limits, a substantial subset already carry
early AD pathology: elevated CSF tau-to-amyloid ratios, amyloid-PET positivity,
and accelerated hippocampal atrophy are detectable at this stage. Identifying
which SCD individuals are on a trajectory to conversion is one of the key
clinical challenges that disease-modifying therapies (DMTs) now demand.

The second stage, \emph{Mild Cognitive Impairment} (MCI), is operationally
defined as objective impairment on cognitive testing---typically one to two
standard deviations below age-adjusted norms---combined with preserved
functional independence in daily activities \citep{albert2011mci,
petersen2004mci}. The amnestic subtype of MCI carries an annual conversion rate
to dementia of approximately 10--15\%, compared with 1--2\% in cognitively
normal individuals. Conventional MCI criteria, however, are susceptible to
false-positive diagnostic errors; \citet{edmonds2015mci} demonstrated that up
to one-third of individuals classified as MCI using standard cut-offs were
cognitively normal on comprehensive neuropsychological testing. This diagnostic
ambiguity makes MCI the hardest boundary for both clinicians and classifiers,
a pattern reflected in this project's per-class results: in Section~\ref{sec:hitl-performance},
MCI shows the lowest AI-only accuracy and the largest gain (+20.3 percentage
points) when the HITL escalation engine routes uncertain cases to human review.

\emph{Alzheimer's dementia} is diagnosed when cognitive deficits are
sufficiently severe to interfere with daily functioning
\citep{mckhann2011diagnosis}. By this stage, widespread synaptic loss, cortical
atrophy, and significant NFT burden are established. The clinical utility of AD
diagnostics therefore pivots increasingly towards the pre-dementia stages, where
intervention is most likely to be effective.

For DMT eligibility, staging precision is critical. Lecanemab was trialled in
patients with early AD (MCI or mild dementia with amyloid confirmation),
achieving a 27\% reduction in clinical decline over 18 months
\citep{vandyck2023lecanemab}. Donanemab similarly targeted early symptomatic
patients, achieving a 35\% slowing of progression in the low/medium tau
subgroup \citep{sims2023donanemab}. Both results underscore that accurate
pre-selection---distinguishing SCD converters, amnestic MCI, and established
AD from non-converters and non-AD conditions---translates directly to treatment
benefit, and that the cost of misclassification is not merely statistical but
clinical.

% -----------------------------------------------------------
\section{Clinical Diagnosis and the Role of AI}
\label{sec:clinical-diagnosis}
% -----------------------------------------------------------

Current clinical diagnosis of AD stages relies on the convergence of evidence
from multiple modalities: standardised cognitive assessments such as the
Mini-Mental State Examination (MMSE), Montreal Cognitive Assessment (MoCA), and
the Clinical Dementia Rating (CDR); structural neuroimaging to assess cortical
atrophy and hippocampal volume; CSF biomarkers for A$\beta$ and tau; and
increasingly, plasma biomarkers for NfL and p-tau217. A definitive ante-mortem
diagnosis still requires biomarker confirmation under most research-grade
protocols.

This multi-modal assessment process is resource-intensive and demands
specialised expertise. The classification boundaries between SCD, MCI, and AD
are not hard thresholds but probabilistic gradients, requiring clinicians to
exercise considerable judgement when interpreting borderline biomarker
profiles. \citet{edmonds2015mci} showed that conventional criteria are
susceptible to false-positive errors, and inter-rater agreement between
specialists can be moderate---introducing both under-diagnosis (missing
patients who would benefit from early treatment) and over-diagnosis (exposing
patients to unnecessary follow-up or DMTs with significant side-effect
profiles such as amyloid-related imaging abnormalities).

Supervised machine learning offers a complementary path. Models trained on
large, curated biomarker datasets can identify multivariate feature
interactions that predict staging labels more reliably than individual
biomarkers in isolation, and can apply these patterns consistently without
fatigue or inter-rater drift \citep{rathore2017review}. Yet applying ML to
clinical decision support introduces well-documented challenges: opaque
predictions erode clinician trust, high-confidence errors propagate
uncorrected, and there is no inherent mechanism for incorporating domain
expertise or patient-specific context that lies outside the feature
space \citep{topol2019highperformance}.

A HITL framework addresses these shortcomings structurally. The ML system
serves as a first-line screener that quantifies its own uncertainty through
multiple complementary signals---entropy, model disagreement, conformal
prediction set size---and escalates cases that exceed a calibrated threshold to
human review. This design preserves the throughput and consistency of
automated classification for clear-cut cases while reserving clinician
authority for the ambiguous boundary cases where human expertise adds the
greatest marginal value. As the results in Chapter~\ref{chap:results}
demonstrate, this selective routing reduces clinical misclassification cost by
62.4\% at a 40\% review budget, with the largest gains concentrated in exactly
the MCI boundary cases where diagnostic uncertainty is highest.

% -----------------------------------------------------------
\section{The ADNI Dataset}
\label{sec:adni}
% -----------------------------------------------------------

The Alzheimer's Disease Neuroimaging Initiative (ADNI) is a longitudinal,
multi-site observational study launched in 2004 with the primary goal of
validating biomarkers for use as clinical trial endpoints
\citep{weiner2017adni}. Now in its fourth phase (ADNI4), the initiative has
enrolled over 2{,}400 participants across more than 60 sites in the United
States and Canada, collecting repeated measurements of MRI volumetrics, amyloid
and tau PET imaging, CSF and plasma biomarkers, genetic markers (including APOE
genotyping), and comprehensive neuropsychological batteries.

The dataset used in this project is constructed by merging four ADNI source
tables on participant identifier (RID) and visit code (VISCODE):

\begin{enumerate}[nosep]
  \item \textbf{ADNIMERGE} --- the canonical longitudinal table providing
    demographics (age, sex, education), CSF biomarkers (TAU, PTAU,
    A$\beta_{1\text{--}42}$), MMSE scores, and diagnostic labels;
  \item \textbf{BLENNOWPLASMANFL} --- plasma neurofilament light chain
    (PLASMA\_NFL);
  \item \textbf{UPENNPLASMA} --- plasma A$\beta_{42}$ and A$\beta_{40}$, from
    which the AB42/40 ratio is computed;
  \item \textbf{BLENNOWPLASMATAU} --- plasma total tau (PLASMATAU).
\end{enumerate}

After merging and dropping records without plasma NfL measurements, the cohort
comprises \textbf{576 patients} with a near-balanced class distribution:
193~SCD (33.5\%), 195~MCI (33.9\%), and 188~AD (32.6\%). The mean age is
75.2~years (SD~6.7), 58.2\% are male, and the mean MMSE score is 26.5
(range 18--30). The MMSE distribution across stages---SCD: $29.1 \pm 1.0$,
MCI: $26.9 \pm 1.8$, AD: $23.3 \pm 2.1$---illustrates the overlapping
cognitive profiles that make three-class separation challenging, particularly
at the SCD--MCI boundary where score distributions overlap substantially.

The raw 13-column feature set is extended by the preprocessing pipeline
(Chapter~\ref{chap:design}) with five engineered ratio features
(e.g.\ TAU/PTAU, ABETA/TAU, PTAU/ABETA), yielding 16 candidate features.
Mutual-information-based feature selection (SelectKBest, $k{=}12$) retains the
12 most informative features, dropping sex (PTGENDER), education (PTEDUCAT),
and the AB42/40 ratio (AB4240) whose marginal information gain falls below the
selection threshold.

Missingness is concentrated in the CSF modality: TAU, PTAU, and ABETA are each
missing for 33.9\% of patients, reflecting the invasiveness of lumbar
puncture; AB4240 is missing for 10.9\%; and PLASMATAU for just 0.5\%. All
remaining features---including PLASMA\_NFL, MMSE, AGE, and the demographic
variables---are fully observed. This missingness pattern is clinically
realistic and mirrors the differential availability of CSF versus plasma
assays in routine practice. The pipeline handles it using distance-weighted
$k$-nearest-neighbour imputation ($k{=}9$), applied within each
cross-validation fold to prevent information leakage from held-out samples
into imputed training values.

Two properties of this cohort are important for interpreting downstream results.
First, the near-equal class balance eliminates prior-probability confounding,
meaning that accuracy differences across classes (Section~\ref{sec:hitl-performance})
reflect genuine classifier difficulty rather than prevalence imbalance.
Second, the 33.9\% CSF missingness means that a substantial fraction of
patients are classified primarily from plasma, cognitive, and demographic
features---a scenario where model uncertainty is expected to be higher and the
HITL escalation engine (Section~\ref{sec:escalation-quality}) is designed to
intervene.


================================================================
  CHAPTER 3: TECHNICAL BACKGROUND AND LITERATURE REVIEW
  Sections 3.5 -- 3.7
================================================================

% -----------------------------------------------------------
\section{Explainable AI}
\label{sec:xai}
% -----------------------------------------------------------

Clinician trust in AI-assisted diagnosis depends not only on prediction accuracy
but on the intelligibility of the reasoning behind each prediction.
Explainable AI (XAI) methods bridge this gap by attributing contributions to
individual input features, enabling a clinician to assess whether the model's
reasoning aligns with established biomarker knowledge before acting on its
recommendation.

SHAP (SHapley Additive exPlanations) \citep{lundberg2017shap} provides a
game-theoretic framework for feature attribution grounded in Shapley values from
cooperative game theory. For a prediction $f(\mathbf{x})$, each feature $i$
receives an attribution $\phi_i$ satisfying:
\begin{equation}
  f(\mathbf{x}) = \phi_0 + \sum_{i=1}^{p} \phi_i,
  \label{eq:shap}
\end{equation}
where $\phi_0 = \mathbb{E}[f(\mathbf{X})]$ is the expected model output across
the dataset and $\sum_i \phi_i$ captures the deviation from this baseline for
the individual patient. The Shapley value for feature $i$ is defined as the
average marginal contribution of that feature across all possible feature
coalitions, which guarantees three properties that are valuable in a clinical
context: \emph{efficiency} (attributions sum exactly to the prediction),
\emph{symmetry} (features with identical contributions receive equal
attribution), and \emph{missingness} (features absent from the input contribute
zero). Together, these ensure that SHAP attributions are both faithful to the
model and internally consistent---a clinician examining two patients with
identical biomarker profiles will see identical explanations.

For tree-based models, SHAP's TreeExplainer computes exact Shapley values in
polynomial time by exploiting the tree structure to enumerate feature
coalitions efficiently \citep{lundberg2017shap}. This computational tractability
is critical for two reasons. First, it enables real-time explanation generation
in the clinical dashboard (Chapter~\ref{chap:design}), where a SHAP waterfall
plot is rendered for each escalated patient at the point of clinician review.
Second, it makes it feasible to compute SHAP values for the full 460-sample OOF
set across all three classes, yielding a $460 \times 12 \times 3$ attribution
tensor that supports both global importance ranking and per-class directional
analysis (Section~\ref{sec:explainability}).

Beyond individual explanations, SHAP attributions serve a second role in this
project as an escalation signal. The \emph{SHAP instability} metric computes the
mean per-feature variance of attribution vectors across the base models for a
given patient. When different models attribute importance to different features
for the same case, the ensemble's prediction lacks a consistent biomarker
grounding even if the predicted class is unanimous. High instability is
therefore a meaningful warning signal that complements the probabilistic
uncertainty measures described in Section~\ref{sec:uq}, and is included as one
of the seven signals in the composite escalation risk score.

% -----------------------------------------------------------
\section{Active Learning and Feedback Loops}
\label{sec:active-learning}
% -----------------------------------------------------------

Active learning is a semi-supervised paradigm in which the model identifies
which unlabelled examples would most improve its performance if labelled, and
selectively queries an oracle---typically a human annotator---for those labels
\citep{settles2009active}. The key insight is that not all labelled corrections
are equally informative: a clinician who confirms a confident, correct
prediction adds no signal, while a correction on a genuinely uncertain boundary
case provides high epistemic value because it resolves ambiguity in precisely
the region of feature space where the model is weakest.

\citet{settles2009active} categorises query strategies into three families.
\emph{Uncertainty sampling} prioritises examples with the highest predictive
entropy or lowest confidence. \emph{Query by committee} selects examples on
which a committee of models disagrees most strongly, measuring disagreement
through vote entropy or KL divergence. \emph{Expected model change} selects
examples whose labels would produce the largest gradient update, directly
optimising for information gain. In an ensemble-based HITL system, query by
committee aligns naturally with the model disagreement signal already computed
for escalation purposes---the cases that the escalation engine routes for
review are, by construction, the cases on which the ensemble's constituent
models disagree most.

In this project, the feedback loop implements a risk-ranked retraining strategy
that extends uncertainty sampling. Clinician corrections are accumulated into a
feedback pool ordered by the composite escalation risk score. After each batch
of corrections, the full ensemble (CatBoost, Random Forest, and neural network)
is retrained on the union of the original training set and the feedback pool.
Each corrected sample is weighted by a function of the clinician's self-reported
confidence on a 1--5 scale: $w = 0.5 + c/5$, where $c$ is the confidence
rating, yielding weights in $[0.7, 1.5]$. This weighting ensures that
high-confidence corrections exert proportionally greater influence on the
retrained model, providing a principled mechanism for handling the annotation
noise that is inherent in simulated (and, eventually, real) clinician feedback.

% -----------------------------------------------------------
\section{Summary and Research Positioning}
\label{sec:positioning}
% -----------------------------------------------------------

The preceding sections characterise a research landscape in which several
active threads---biomarker-based AD classification, ensemble learning with
calibration, HITL systems, conformal prediction, explainable AI, and active
learning---have largely developed in isolation. Existing ML approaches to AD
staging predominantly optimise predictive accuracy on held-out test sets without
modelling the downstream clinical workflow in which predictions will be consumed.
HITL frameworks in the medical AI literature have been studied primarily in
radiology and image annotation contexts \citep{rajpurkar2022ai}, where the
modality and annotation process differ qualitatively from tabular biomarker data.
Conformal prediction has been applied to medical classification tasks but has
not, to the authors' knowledge, been integrated into a multi-signal escalation
engine for AD staging. Similarly, SHAP-based explainability and active learning
feedback loops are well-established individually, but their combination within a
single end-to-end AD staging pipeline has not been reported.

This project occupies a specific gap at the intersection of these threads. Its
contributions are as follows:

\begin{enumerate}[nosep]
  \item Stacked generalisation with probability calibration for three-class AD
    staging on ADNI data, evaluated via leak-free out-of-fold cross-validation
    with a separate held-out test set.
  \item A multi-signal escalation engine combining seven distinct uncertainty and
    data-quality indicators under a three-tier policy, with formal ablation of
    each signal's contribution.
  \item Conformal prediction sets (APS) providing finite-sample coverage
    guarantees that are directly communicable to clinicians as set-valued
    diagnoses.
  \item An interactive clinical dashboard with real-time SHAP explanations, a
    structured clinician feedback mechanism, and an active-learning retraining
    loop.
  \item A cost--benefit analysis under a clinically motivated asymmetric error
    matrix and a fairness audit across demographic subgroups
    \citep{mehrabi2021survey}.
\end{enumerate}

Together, these components constitute a more complete HITL framework for AD
staging than has been reported in the existing literature: one that spans from
raw biomarker data through uncertainty-aware classification, selective
escalation, explainable review, and closed-loop model improvement.

